%% ======================================================================
\section{Self-bounding learning}
\label{chap:selfbounding}
%% ======================================================================

\begin{objectives}
\guiding{If a bound holds for every posterior, why not choose the posterior that
makes it smallest, and what does this change to the certificate?}
Up to now, PAC-Bayesian bounds were used after learning, to certify a posterior
produced by some algorithm. This section reverses the roles: since the bounds
hold for all posteriors, we can \emph{minimize} them, and obtain a predictor
together with its own risk certificate. We explain why this is legitimate, what
may and may not depend on the data, how this approach compares with the
classical hold-out certificate, and which optimization tools are used by all
self-bounding algorithms. Each point is illustrated by running the algorithms on
random forests.
\end{objectives}

The bounds of Sections~\ref{chap:bounds} to~\ref{chap:mv} have a remarkable
property: they hold \emph{simultaneously for all posteriors}. So far we have
used them to evaluate a posterior produced by another algorithm, such as the
uniform weights of a random forest; in our running example, the uniform
forest received a tandem certificate of $0.387$. But nothing prevents us from choosing the
posterior \emph{because} its bound is small. Instead of minimizing an empirical
risk and then trying to certify the result, we can directly minimize the bound:
the algorithm then outputs a predictor \emph{and} a guarantee on its risk. This
idea goes back to \citet{freund1998self} and \citet{langford2003microchoice}, who
coined the term \emph{self-bounding learning algorithm}. PAC-Bayes turned it
into a general and practical methodology
\citep{germain2009pac,germain2015risk,thiemann2017strongly,masegosa2020second,viallard2021self,perez2021tighter}.

%% ----------------------------------------------------------------------
\subsection{From empirical risk minimization to bound minimization}

To see what self-bounding learning changes, let us first recall how we usually
learn and evaluate a model. In the first part of the course, learning followed
a three-step pipeline. A model is learned by minimizing a regularized empirical risk on a training
sample, its hyperparameters are tuned by cross-validation, and it is finally
evaluated on a test sample that was never used before. The test error is an
\emph{estimate} of the true risk, which a concentration inequality turns into a
\emph{guarantee} (the test-set bound of \cref{sec:holdout}). This last step has a
cost: the examples of the test sample do not contribute to learning.

A self-bounding algorithm merges the three steps into one. Suppose that we have a
PAC-Bayesian bound of the form
\[
\P_{S\sim\D^m}\Big(\forall Q\in\M(\Hcal),\ \ \RD(Q)\leq\mathcal B(Q,S,\delta)\Big)\geq1-\delta,
\]
where $\RD(Q)$ denotes the risk we want to control (the Gibbs risk, or the
risk of the vote) and $\mathcal B(Q,S,\delta)$ is a quantity that can be computed
from the sample. The algorithm learns
\begin{equation}\label{eq:sb_principle}
Q_S\in\argmin_{Q\in\mathcal Q}\ \mathcal B(Q,S,\delta),
\qquad\text{and returns the certificate}\qquad \mathcal B(Q_S,S,\delta),
\end{equation}
where $\mathcal Q\subseteq\M(\Hcal)$ is the family of posteriors in which we
search, for instance all distributions on a finite set of voters, or Gaussian
distributions on the weights of a model.

\begin{definition}{Self-bounding learning algorithm}{selfbounding}
A \emph{self-bounding learning algorithm} is a procedure which, given a sample
$S\sim\D^m$ and a confidence level $\delta\in(0,1]$, returns a predictor $f_S$ and a
number $\mathcal B_S\in[0,1]$, both computed from $S$ only, such that
\[
\P_{S\sim\D^m}\Big(\RD(f_S)\leq\mathcal B_S\Big)\geq1-\delta.
\]
The number $\mathcal B_S$ is called the \emph{risk certificate} of $f_S$.
\end{definition}

In plain words, a self-bounding algorithm hands us two things computed from the
same data: a predictor, and a number that we can trust, with high probability,
as an upper bound on its true risk. Many algorithms of this course fit this
definition: the Gibbs posterior
$Q_\lambda$ of \cref{sec:gibbs_posterior} with its Seeger certificate, the
forest weighted by \cref{algo:fo} with its first-order certificate, or PBGD with
Catoni's bound (\cref{sec:linear}). A random forest with uniform weights,
returned together with its out-of-bag tandem certificate, is also
self-bounding in the sense of the definition; the certificate is simply not used
to learn. One question must be settled before we go further: if the posterior is
chosen by looking at its own bound, is that bound still trustworthy?

%% ----------------------------------------------------------------------
\subsection{Why the certificate remains valid}

The following result is almost trivial, but it is the reason why the whole
approach works, so it deserves to be stated carefully.

\begin{theorem}{Validity of self-bounding certificates}{sb_valid}
Let $\mathcal B$ be such that $\P_S\big(\forall Q\in\mathcal Q,\ \RD(Q)\leq\mathcal B(Q,S,\delta)\big)\geq1-\delta$.
Let $S\mapsto Q_S\in\mathcal Q$ be \emph{any} learning rule (deterministic, or using a
source of randomness independent of the data). Then
\[
\P_S\Big(\RD(Q_S)\leq\mathcal B(Q_S,S,\delta)\Big)\geq1-\delta .
\]
In particular this holds for $Q_S$ minimizing $\mathcal B(\cdot,S,\delta)$ or any
other objective, and for any stopping time of an iterative optimizer.
\end{theorem}
\begin{proof}
Let $E=\{S:\forall Q\in\mathcal Q,\ \RD(Q)\leq\mathcal B(Q,S,\delta)\}$. For every
$S\in E$, the inequality holds in particular for the posterior $Q_S$. Hence
$\{S:\RD(Q_S)\leq\mathcal B(Q_S,S,\delta)\}\supseteq E$ and its probability is at
least $\P(E)\geq1-\delta$. If $Q_S$ also depends on an independent random seed $\xi$,
apply the argument conditionally on $\xi$ and integrate.
\end{proof}

In words: since the bad event does not depend on which posterior we pick, no
way of picking one, however clever or data-driven, can make us fall into it.
The key is the quantifier ``for all $Q$'' inside the probability. A bound which
only holds for a \emph{fixed} predictor, such as Hoeffding's inequality for a
fixed $h$ or a test-set bound, does not remain valid if the predictor is chosen
by looking at the data used in the bound.

We saw in \cref{fig:uniformity} how
badly such a bound fails: after selecting the best of $1000$ random classifiers,
the naive Hoeffding bound was wrong in $99.8\%$ of the runs. The Occam bound
remained valid because it pays $\ln\frac{1}{P(h)}=\ln1000$, which is exactly the
price of having \emph{selected} $h$ with the data. PAC-Bayes extends this
accounting to any soft selection: the price of choosing a posterior $Q$ with the
data is $\KL(Q\|P)$. Because the price is already included in the bound, we are
allowed to choose $Q$ in any way we like, and in particular by minimizing the
bound. This freedom concerns the posterior only, and the next subsection makes
precise where it stops.

%% ----------------------------------------------------------------------
\subsection{What may depend on the data?}\label{sec:sb_rules}

Theorem~\ref{thm:sb_valid} allows the posterior to depend arbitrarily on $S$.
Everything else in the bound --- the prior, the confidence level $\delta$, the form
of the bound and its internal parameters --- must be fixed before seeing the
sample used in the bound. When we nevertheless want to choose one of these
ingredients with the data, we pay for it with a union bound, as the following
proposition shows.

\begin{proposition}{Combining several bounds}{union}
Let $\mathcal B_1,\dots,\mathcal B_K$ be $K$ bounds (they may use different priors,
different parameters, different forms, or even different sets of voters),
each valid in the sense that for every $\delta'$,
$\P_S\big(\forall Q,\ \RD(Q)\leq\mathcal B_k(Q,S,\delta')\big)\geq1-\delta'$.
Let $\pi_1,\dots,\pi_K>0$ with $\sum_k\pi_k=1$. Then
\[
\P_S\Big(\forall k,\ \forall Q,\ \ \RD(Q)\leq\mathcal B_k(Q,S,\pi_k\delta)\Big)\geq1-\delta.
\]
Consequently, $\min_k\mathcal B_k(Q_S,S,\pi_k\delta)$ is a valid certificate for any
learning rule $Q_S$, and the index $k$ may be chosen with the data.
\end{proposition}
\begin{proof}
The probability that at least one of the $K$ events fails is at most
$\sum_k\pi_k\delta=\delta$ (union bound).
\end{proof}

In other words, we may try several bounds and keep the best one, provided we
share the confidence budget $\delta$ among them in advance. With
$\pi_k=\frac1K$, choosing among $K$ candidates amounts to replacing $\delta$ by
$\delta/K$, that is to adding $\ln K$ to the numerator of the complexity term.
This is a very small price for a moderate $K$: choosing among ten candidates
costs $2.3$ nats, the same as dividing $\delta$ by ten. With a countable family,
one can take $\pi_k=\frac{6}{\pi^2k^2}$, which costs $2\ln k+\ln\frac{\pi^2}6$ for the
$k$-th candidate.

The line between what is free and what must be paid for is subtle, so the
table below gathers the most common cases. The underlying rule is simple to
state. A parameter is free when it only indexes a family of
\emph{posteriors}, all compared with the same prior in the same bound: the
parameter $C$ of PBGD, for instance, selects a weight vector $\mathbf w$, hence a
posterior $Q_{\mathbf w}$, and Seeger's bound holds for all of them at once. A
parameter is not free when it changes the \emph{probabilistic event} on which
the bound holds: the parameter $C$ of Catoni's bound changes the function
$\Delta$ of Theorem~\ref{thm:general}, and each value gives a different event.

\begin{center}
\renewcommand{\arraystretch}{1.3}
\begin{tabular}{p{4.6cm}p{4.2cm}p{4.2cm}}
\toprule
\textbf{Free} (no correction) & \textbf{Union bound} ($\delta\to\delta/K$) & \textbf{Forbidden} \\
\midrule
\small the posterior $Q$ and the algorithm producing it; the temperature of a Gibbs posterior; $\lambda$ in the PAC-Bayes-$\lambda$ bound; $C$ in PBGD; the learning rate, initialization and number of iterations (early stopping on the bound)
&
\small $C$ in Catoni's bound; $\lambda$ in the Hoeffding-type bound; the variance of a Gaussian prior chosen on a grid; the choice among first-order / tandem / C-bound; the choice among $K$ ensembles built with different hyperparameters
&
\small a prior learned on the examples used in the bound; choosing $\delta$ after seeing the bound; repeating the experiment with many seeds and reporting the best certificate (unless all of them are accounted for) \\
\bottomrule
\end{tabular}
\end{center}

\paragraph{Self-bounding model selection.} Because a choice among $K$ candidates
costs only $\ln K$, Proposition~\ref{prop:union} offers a principled alternative
to cross-validation. We build $K$ candidate models, for
instance forests with different tree depths, learn a posterior for each of them
by minimizing its bound computed with $\delta/K$, and keep the candidate with the
smallest certificate. The selected model comes with a valid certificate, and no
example was set aside for validation.

\begin{example}[Choosing the depth of the trees with the certificate]
On \texttt{digits}, we build six forests of $100$ trees, with maximal depths $1$,
$2$, $3$, $5$, $8$ and unlimited, each tree being trained on a bootstrap sample
of half the training set. For each forest, the posterior minimizes the
out-of-bag tandem bound of \cref{sec:oob}, computed with $\delta/6$.
\Cref{fig:sb_selection} reports the certificates and the test risks over five
random splits. The certificate decreases steadily with the depth, from $0.96$
for stumps to $0.375$ for unpruned trees, and so does the test risk of the vote,
from $0.19$ to $0.031$. The certificate thus selects the unpruned trees, which
is also the best choice for the test risk, and the price of the selection
($\ln6\approx1.8$ nats) is invisible at this scale. (These values are averaged over five
stratified splits, which explains the small differences with the single split
used in \cref{chap:algos}.) In general, the certificate remains a pessimistic
proxy for the test risk: when two candidates are close, it may select a slightly
sub-optimal one.
\end{example}

\begin{figure}[t]
\centering
\includegraphics[width=0.66\linewidth]{selfbounding_selection}
\caption{Self-bounding model selection. Six forests (100 trees, bootstrap samples
of size $m/2$) with increasing tree depth are learned on \texttt{digits}; for each
of them the posterior minimizes the tandem bound and the certificate is computed
with $\delta/6$ (Proposition~\ref{prop:union}). The certificate decreases with the depth
like the test risk; here it selects the model with the best test risk (means and standard deviations over 5 splits).}
\label{fig:sb_selection}
\end{figure}

%% ----------------------------------------------------------------------
\subsection{Self-bounding versus hold-out}\label{sec:holdout}

A self-bounding certificate is only interesting if it compares well with the
simplest alternative. This natural competitor is the \emph{test-set
bound} \citep{langford2005tutorial}, which certifies a predictor with an
independent sample.

\begin{proposition}{Test-set bound}{testset}
Let $f$ be a predictor learned on $S_1$ and $T\sim\D^n$ a sample independent of $S_1$.
Then, with probability at least $1-\delta$ over $T$,
$\RD(f)\leq\klup\big(\widehat R_T(f),\frac{\ln(1/\delta)}{n}\big)$.
\end{proposition}
\begin{proof}
Conditionally on $S_1$, $f$ is fixed and $n\widehat R_T(f)\sim\mathrm{Bin}(n,\RD(f))$;
apply Chernoff's bound (Lemma~\ref{lem:chernoff}). Equivalently, it is the Occam
bound with a single hypothesis ($P=\delta_f$, $\KL=0$).
\end{proof}

The test-set bound certifies \emph{any} predictor, in particular the majority
vote itself, with no factor $2$ or $4$ and no KL term, so it is very tight. Its
drawback is that the $n$ test examples do not contribute to learning.
Self-bounding learning makes the opposite trade-off: all the data serve both to
learn and to certify, but the certificate relies on a PAC-Bayesian bound, which
is looser. Which trade-off wins is an empirical question, which the next
example examines.

\begin{example}[Which approach gives the better vote, and the better certificate?]
We compare the two approaches on a simulated problem, for sample sizes $m$ from
$200$ to $4000$ (\cref{fig:sb_holdout}). The self-bounding approach builds a
forest of $100$ trees on the $m$ examples and minimizes its out-of-bag tandem
bound; the hold-out approach builds the same forest on $70\%$ of the examples
and computes the test-set bound of its vote on the remaining $30\%$. As
expected, using all the data gives a better vote for every $m$: with $m=1000$,
its test risk is $0.099$ against $0.110$. But the hold-out certificate is much
tighter: $0.17$ against $0.63$ for $m=1000$, and $0.10$ against $0.39$ for
$m=4000$. For majority votes, the factor $4$ of the tandem bound and the
second-order KL term cost more than the examples withheld by the hold-out.
\end{example}

\begin{figure}[t]
\centering
\includegraphics[width=\linewidth]{selfbounding_vs_holdout}
\caption{Self-bounding (forest learned on the $m$ examples, posterior minimizing the
OOB tandem bound, TND certificate) versus hold-out (forest learned on $70\%$ of
the examples, test-set bound on the remaining $30\%$), on a simulated problem,
as a function of $m$ (average over 3 repetitions). Using all the data gives a
better vote (left), but for majority votes the hold-out certificate remains much
tighter (right).}
\label{fig:sb_holdout}
\end{figure}

This experiment should be read honestly. If the only goal is to report a tight
guarantee on a fixed vote, the hold-out is hard to beat. The value of
self-bounding learning lies elsewhere: it uses all the data, it provides a
learning objective with a guarantee, and it makes model selection almost free.

The gap in tightness is also specific to majority votes, where the passage from
the Gibbs risk to the vote costs a constant factor; for stochastic neural
networks, self-bounding certificates with data-dependent priors are tight
enough to be informative, for instance on MNIST \citep{perez2021tighter}. The
following table summarizes the comparison.

\begin{center}
\renewcommand{\arraystretch}{1.3}
\begin{tabular}{p{3.2cm}p{4.8cm}p{4.8cm}}
\toprule
& Hold-out (test-set bound) & Self-bounding (PAC-Bayes) \\
\midrule
Data used to learn & part of the sample & the whole sample \\
Predictor certified & any (vote, deep net\dots) & a posterior (Gibbs, vote via Section~\ref{chap:mv}) \\
Tightness & excellent: a single kl inversion, no KL term & depends on the KL and on the Gibbs-to-vote relation \\
Guides learning? & no & yes: the certificate is the objective \\
Model selection & needs another split & free or $\ln K$ (Proposition~\ref{prop:union}) \\
\bottomrule
\end{tabular}
\end{center}

\paragraph{Best of both worlds: data-dependent priors.} We do not have to choose
between the two approaches: they can be combined. We split $S=S_1\cup S_2$, use $S_1$ to learn a \emph{prior} $P_{S_1}$
(a Gaussian centred on a network trained on $S_1$, or voters built on $S_1$),
and compute the PAC-Bayesian bound on $S_2$.

\begin{proposition}{Self-bounding with a data-dependent prior}{ddprior}
Let $P_{S_1}$ be a prior learned on $S_1$ and $\mathcal B$ a PAC-Bayesian bound
computed on $S_2\sim\D^{m_2}$, independent of $S_1$. Then, with probability at least
$1-\delta$ over $(S_1,S_2)$, for all posteriors $Q$ --- including posteriors learned on the whole
sample $S_1\cup S_2$ --- we have $\RD(Q)\leq\mathcal B(Q,S_2,P_{S_1},\delta)$.
\end{proposition}
\begin{proof}
Conditionally on $S_1$, the prior is fixed and $S_2$ is an i.i.d.\ sample, so the
bound holds with probability at least $1-\delta$ over $S_2$, for all $Q$ simultaneously.
Integrate over $S_1$.
\end{proof}

Put simply, the prior may look at $S_1$ as long as the bound only looks at
$S_2$. Only the \emph{empirical term} of the bound must be computed on $S_2$; the
posterior itself may use all the data. This is the strategy behind the tightest
certificates for neural networks \citep{dziugaite2017computing,perez2021tighter}
and SVMs \citep{ambroladze2006tighter,parrado2012pac}, and we measured its
trade-off in \cref{fig:ddprior}.

The two approaches of this subsection appear as the two ends of this spectrum.
The test-set bound is the extreme case where
the posterior is a Dirac on the prior ($\KL=0$), and a pure self-bounding
algorithm is the other extreme ($S_1=\emptyset$). The out-of-bag construction of
\cref{sec:oob} can be seen as a clever way to obtain a data-dependent prior
without splitting the sample once and for all.

%% ----------------------------------------------------------------------
\subsection{Anatomy of a self-bounding algorithm}

We now know what is allowed; it remains to turn these rules into algorithms.
Fortunately, all the self-bounding algorithms of this course follow the same
template, summarized in \cref{fig:sb_pipeline} and \cref{algo:sb_template}. The
first decision is how to organize the data: which examples build the voters or the
prior, and which ones are used in the bound. The prior is then fixed,
independently of the examples of the bound. The third step chooses an
objective, which may be the bound itself or a relaxation that is easier to
optimize, and the fourth step minimizes it with one of the tools of
\cref{sec:sb_optim}. Finally, the tightest valid bound is computed for the
posterior that was found.

\begin{figure}[t]
\centering
\begin{tikzpicture}[node distance=0.5cm and 0.32cm,
   box/.style={draw,rounded corners,align=center,minimum height=1.3cm,text width=2.05cm,font=\footnotesize,fill=#1!6,draw=#1},
   arr/.style={-{Stealth[length=2.2mm]},thick}]
\node[box=pbteal] (a) {\textbf{1. Voters}\\ and data\\ organization\\ (split, OOB)};
\node[box=pbgrey,right=of a] (b) {\textbf{2. Prior} $P$\\ fixed before\\ the bound sample};
\node[box=pbviolet,right=of b] (c) {\textbf{3. Objective}\\ relaxed bound\\ $\widetilde{\mathcal B}(Q)$};
\node[box=pbblue,right=of c] (d) {\textbf{4. Optimize}\\ $Q_S\approx\argmin\widetilde{\mathcal B}$};
\node[box=pbred,right=of d] (e) {\textbf{5. Certify}\\ tight kl form\\ at $Q_S$};
\draw[arr] (a) -- (b); \draw[arr] (b) -- (c); \draw[arr] (c) -- (d); \draw[arr] (d) -- (e);
\draw[arr,dashed,pbgrey] (e.south) -- ++(0,-0.5) -| node[pos=0.25,below,font=\small]{model selection with $\delta/K$ (Prop.~\ref{prop:union})} (a.south);
\end{tikzpicture}
\caption{The five ingredients of a self-bounding algorithm.}
\label{fig:sb_pipeline}
\end{figure}

\begin{algorithm}[h]
\caption{Generic self-bounding learning}\label{algo:sb_template}
\KwIn{sample $S$, confidence $\delta$, family of posteriors $\{Q_\theta\}$}
Organize the data: which examples build the voters / the prior, which ones are used in the bound\;
Fix the prior $P$ (independent of the examples used in the bound)\;
Choose a \emph{relaxed} objective $\widetilde{\mathcal B}(\theta)\geq$ (or $\approx$) the bound, easy to optimize\;
$\theta_S\leftarrow$ approximately minimize $\widetilde{\mathcal B}(\theta)$ (closed form, alternating or gradient-based)\;
Compute the \emph{tight} certificate $\mathcal B(Q_{\theta_S},S,\delta)$ with the kl inversion (Monte Carlo if needed)\;
\KwOut{$Q_{\theta_S}$ and its certificate}
\end{algorithm}

Two features of this template are worth stressing. First, the objective need
not be the certificate. By Theorem~\ref{thm:sb_valid}, the certificate is valid
whatever the way $Q_S$ was obtained, so we may optimize a convenient relaxation
(McAllester's bound, the PAC-Bayes-$\lambda$ bound, a surrogate loss) and report
the tightest valid bound at the end; if we report the minimum of several bounds,
Proposition~\ref{prop:union} applies. Second, the optimizer need not find the
global minimum: every posterior is certified, and a poor optimizer only yields
a looser certificate, never a wrong one. Step~4 is therefore a matter of
efficiency rather than validity, and the next subsection reviews the tools
available for it.

%% ----------------------------------------------------------------------
\subsection{Optimization toolbox}\label{sec:sb_optim}

The right optimizer depends on the shape of the objective. We go from the
easiest case, where the minimizer is explicit, to the most general one, where
we differentiate through the kl inversion itself.

\subsubsection{Closed form: linear objectives and the Gibbs posterior}

When the objective is linear in the empirical term and in the KL,
$\RS(\GQ)+\frac{\KL(Q\|P)}{\lambda m}$, as for Hoeffding-type or PAC-Bayes-$\lambda$
bounds, Proposition~\ref{prop:gibbs_post} gives the minimizer in closed form,
$Q_\lambda(h)\propto P(h)e^{-\lambda m\RS(h)}$. The only remaining parameter is the
temperature, which is free since it indexes posteriors: we can scan a grid of
values of $\lambda$ and keep the best certificate, as in \cref{fig:temperature}.

\subsubsection{Alternating minimization}

When the temperature is not scanned but optimized, a closed form is available
in each variable separately. The PAC-Bayes-$\lambda$ bound \eqref{eq:lambda} is minimized in $Q$ by the Gibbs
posterior when $\lambda$ is fixed, and in $\lambda$ by the closed form
\eqref{eq:lambda_star} when $Q$ is fixed. Alternating the two steps
(\cref{algo:fo}) converges to the global minimum under the conditions of
\citet{thiemann2017strongly}; we will see in \cref{tab:fo_iterations} that a few
iterations suffice in practice.

\subsubsection{Gradient-based minimization}

For richer objectives, such as second-order bounds, C-bounds or bounds on
neural networks, we use gradient descent on a parametrization of the posterior.
For a finite set of voters, the simplest parametrization is the softmax
$\rho=\mathrm{softmax}(\theta)$, $\rho_t=e^{\theta_t}/\sum_se^{\theta_s}$, which
removes the simplex constraint. If $g=\nabla_\rho\mathcal B$ is the gradient
with respect to the weights, the chain rule gives
\begin{equation}\label{eq:softmax_grad}
\nabla_\theta\mathcal B=\rho\odot\big(g-(\rho^\top g)\mathbf 1\big),
\qquad\text{with, for the KL term,}\qquad
\nabla_\rho\KL(\rho\|\pi)=\ln\frac{\rho}{\pi}+\mathbf 1 .
\end{equation}
For Gaussian posteriors $Q=\mathcal N(\boldsymbol\mu,\mathrm{diag}(\boldsymbol\sigma^2))$ on the
parameters $\mathbf w$ of a model, the KL to a Gaussian prior is explicit, and
the expected loss is differentiated with the \emph{reparametrization trick}:
writing $\mathbf w=\boldsymbol\mu+\boldsymbol\sigma\odot\boldsymbol\varepsilon$ with
$\boldsymbol\varepsilon\sim\mathcal N(0,I)$, we have
$\nabla_{\boldsymbol\mu,\boldsymbol\sigma}\E_{\mathbf w\sim Q}\ell(\mathbf w)=\E_{\boldsymbol\varepsilon}\nabla_{\boldsymbol\mu,\boldsymbol\sigma}\ell(\boldsymbol\mu+\boldsymbol\sigma\odot\boldsymbol\varepsilon)$,
which is estimated with a few samples of $\boldsymbol\varepsilon$
\citep{dziugaite2017computing,perez2021tighter}.

In both cases a difficulty remains: the 0-1 loss is not differentiable. Either the Gibbs risk has a smooth
closed form, as the probit loss of linear classifiers (\cref{eq:probit}), or the
0-1 loss is replaced during learning by a bounded surrogate such as a clipped
cross-entropy. The final certificate is always computed with the 0-1 loss.

\subsubsection{Differentiating through the kl inversion}

The tightest bounds are expressed with $\klup$, which has no closed form. It
can nevertheless be differentiated, which allows us to minimize Seeger-type
certificates directly \citep{reeb2018learning,viallard2021self}.

\begin{proposition}{Derivatives of the kl inversion}{klgrad}
Let $q\in(0,1)$, $\varepsilon>0$ and $p=\klup(q,\varepsilon)$ (or $p=\kllow(q,\varepsilon)$), so that
$\kl(q\|p)=\varepsilon$ and $p\neq q$. Then
\[
\frac{\partial p}{\partial q}=\frac{p(1-p)}{p-q}\,\ln\frac{p(1-q)}{q(1-p)},
\qquad
\frac{\partial p}{\partial\varepsilon}=\frac{p(1-p)}{p-q}.
\]
\end{proposition}
\begin{proof}
Let $F(q,p)=\kl(q\|p)-\varepsilon$. We have $F_q=\partial_q\kl(q\|p)=\ln\frac{q(1-p)}{p(1-q)}$
and $F_p=\partial_p\kl(q\|p)=\frac{p-q}{p(1-p)}\neq0$. By the implicit function theorem,
$p$ is a differentiable function of $(q,\varepsilon)$ with $\frac{\partial p}{\partial q}=-\frac{F_q}{F_p}$
and $\frac{\partial p}{\partial\varepsilon}=\frac{1}{F_p}$.
\end{proof}

So, although $\klup$ has no closed form, its gradient only requires the value
$p$ that we compute anyway. As expected, $\frac{\partial p}{\partial\varepsilon}$
is positive for the upper inversion and negative for the lower one. These
derivatives also tell us how
sensitive a certificate is to its ingredients. At $\hat q=0.1$ and $\varepsilon=0.05$,
where $p\approx0.220$, we find $\frac{\partial p}{\partial\varepsilon}\approx1.43$: with $m=1000$
examples, one more nat of KL increases $\varepsilon$ by $0.001$ and the certificate by
about $0.0014$.

For the tandem certificate of a finite ensemble,
$\mathcal B_{\mathrm{TND}}(\rho)=4\,\klup\big(\rho^\top\widehat{\mathbf L}\rho,\,\varepsilon(\rho)\big)$ with
$\varepsilon(\rho)=\frac{2\KL(\rho\|\pi)+\ln(2\sqrt{n_0}/\delta)}{n_0}$, where $n_0$ is the number of examples on which the tandem losses are computed, the chain rule gives
\begin{equation}\label{eq:grad_tnd_cert}
\nabla_\rho\mathcal B_{\mathrm{TND}}=4\Big[\frac{\partial p}{\partial q}\,2\widehat{\mathbf L}\rho+\frac{\partial p}{\partial\varepsilon}\,\frac2n\Big(\ln\frac\rho\pi+\mathbf 1\Big)\Big],
\end{equation}
to be combined with \eqref{eq:softmax_grad}. The function
\texttt{certificate\_and\_grad} of \texttt{pbtools.py} implements this formula.

\begin{example}[Watching a self-bounding algorithm learn]
\Cref{fig:sb_dynamics} runs this gradient descent (Adam, $400$ iterations, softmax
parametrization, uniform initialization) on the out-of-bag first-order and
tandem certificates of a forest of $100$ trees on \texttt{digits}, and records the
certificate, the test risk of the vote and the effective number of trees
$1/\sum_t\rho_t^2$ along the iterations. The first-order certificate decreases
from $0.46$ to $0.39$, but it achieves this by concentrating the posterior on
fewer and fewer trees: after $200$ iterations, fewer than two trees carry the
whole weight, and the test risk of the vote jumps from $0.044$ to $0.15$. The
tandem certificate decreases from $0.39$ to $0.36$ while keeping an effective
number of about $48$ trees, and the vote remains as accurate as the uniform
forest. This is the phenomenon announced in \cref{chap:mv} and analysed in
\cref{sec:algo_fo}: an objective that ignores the correlations between voters
destroys the diversity that makes the vote work.
\end{example}

\begin{figure}[t]
\centering
\includegraphics[width=\linewidth]{selfbounding_dynamics}
\caption{Direct minimization (Adam, softmax parametrization, gradient of
Proposition~\ref{prop:klgrad}) of the OOB first-order and tandem certificates of a forest of
100 trees on \texttt{digits}. Left: value of the certificate; middle: test risk
of the vote; right: effective number of trees.}
\label{fig:sb_dynamics}
\end{figure}

\subsubsection{Relaxations of Seeger's bound used as objectives}

Differentiating through $\klup$ is not always convenient. When a differentiable
closed form is preferred, one of the following upper
bounds on $\klup(\hat q,\varepsilon)$ can be minimized instead \citep{perez2021tighter}.
All of them follow from the Pinsker-type inequalities of Lemma~\ref{lem:pinsker}.
The PAC-Bayes-$\lambda$ bound, optimized in $\lambda$, coincides with the
``quadratic'' bound, but its linear form is more convenient for alternating
minimization.

\begin{center}
\renewcommand{\arraystretch}{1.5}
\begin{tabular}{lll}
\toprule
Name & Objective & Comment \\
\midrule
McAllester (``classic'') & $\hat q+\sqrt{\varepsilon/2}$ & simplest, loose when $\hat q$ is small \\
PAC-Bayes-$\lambda$ & $\frac{\hat q}{1-\lambda/2}+\frac{\varepsilon}{\lambda(1-\lambda/2)}$ & linear; quadratic at $\lambda^\star$ \\
``Quadratic'' & $\big(\sqrt{\hat q+\varepsilon/2}+\sqrt{\varepsilon/2}\big)^2$ & from $(p-\hat q)^2\leq2p\varepsilon$ \\
Seeger (kl) & $\klup(\hat q,\varepsilon)$ & tightest; gradient by Prop.~\ref{prop:klgrad} \\
\bottomrule
\end{tabular}
\end{center}

%% ----------------------------------------------------------------------
\subsection{Pitfalls and good practice}

Self-bounding learning is simple, but its guarantees are only as good as the
care with which the rules of \cref{sec:sb_rules} are applied. Three mistakes are
particularly common.

\begin{pitfall}[: a prior that has seen the data]
Building the voters on the whole sample and then using a uniform prior over them,
with empirical risks computed on the same sample, is \emph{not} valid: the
voters, hence the prior, depend on the examples of the bound. Use a split, or the
out-of-bag construction of \cref{sec:oob}.
\end{pitfall}

\begin{pitfall}[: hidden selections]
Trying several bounds, several priors or several random seeds and reporting the
best certificate is a selection with the data, which must be accounted for with
Proposition~\ref{prop:union}.
\end{pitfall}

\begin{pitfall}[: certifying the wrong predictor]
A bound on the Gibbs risk is not a bound on the risk of the vote (use
\cref{chap:mv}), nor on the risk of a single network drawn from the posterior
(use disintegrated bounds). State precisely what is certified.
\end{pitfall}

Turning these warnings around, we obtain a short checklist. In practice, a sound self-bounding experiment starts by deciding what is
certified: a Gibbs classifier, a majority vote or a deterministic model. The
prior is then fixed before looking at the examples used in the bound, and every
choice made with the data is accounted for, at the price of $\ln K$ for $K$
candidates. A convenient relaxation can be optimized, but the reported
certificate should be the tight kl form. For majority votes of diverse voters,
second-order bounds should be preferred. Finally, a certificate should always
be reported with its ingredients (empirical term, KL, $m$, $\delta$) and next to
the test error, since a bound is a guarantee and not an estimate.

\begin{keypoints}
A PAC-Bayesian bound holds for all posteriors: minimizing it is legitimate and
yields a predictor together with its risk certificate (Theorem~\ref{thm:sb_valid}).
The posterior, and anything that only indexes posteriors, may depend freely on the
data; the prior, $\delta$ and the parameters of the bound may not, and a choice
among $K$ candidates costs $\ln K$. Hold-out certificates are tighter for majority
votes, but self-bounding learning uses all the data, makes the certificate a
learning objective and makes model selection almost free; data-dependent priors
interpolate between the two. The generic recipe is always the same: organize the
data, fix the prior, optimize a relaxed bound (closed form, alternating
minimization, or gradient through the kl inversion), and certify with the kl form.
\end{keypoints}

\begin{quickchecks}
\item Why does minimizing a PAC-Bayesian bound not invalidate the certificate of the minimizer?
\item Is the temperature of a Gibbs posterior a free parameter? And the parameter $C$ of Catoni's bound?
\item What is the price of choosing the best of $K$ candidate models with their certificates?
\item Give one advantage of the hold-out certificate and one advantage of the self-bounding certificate.
\end{quickchecks}

\begin{exercises}
\begin{exercise}
Reproduce the experiment of \cref{fig:uniformity} (draw $1000$ random classifiers
and $m=100$ labels). How often is the ``fixed-$h$'' Hoeffding bound violated?
And the Occam bound? What happens to the Occam bound with $10^6$ classifiers?
\end{exercise}
\begin{exercise}
Show that the value of $\lambda$ minimizing the Hoeffding-type bound
\eqref{eq:catoni_hoeffding} depends on $\KL(Q\|P)$, and explain why using it
directly is forbidden, whereas using $\lambda^\star$ in the PAC-Bayes-$\lambda$ bound is allowed.
\end{exercise}
\begin{exercise}
Using Proposition~\ref{prop:klgrad}, compute $\partial p/\partial q$ at
$\hat q=0.1$, $\varepsilon=0.05$ (with $p\approx0.220$, see \cref{fig:klinv}). Which ingredient of
the certificate is it more profitable to decrease: the empirical risk by $0.01$,
or the KL by $5$ nats with $m=1000$?
\end{exercise}
\begin{exercise}
$(\star)$ Derive the ``quadratic'' relaxation of the table above from the refined
Pinsker inequality. Writing $a=\sqrt{\varepsilon/2}$, show that it is smaller than
McAllester's bound iff $2a+2\sqrt{\hat q+a^2}\leq1$; in particular, for small
$\varepsilon$ it is the case as soon as $\hat q<\frac14$.
\end{exercise}
\begin{exercise}
\textbf{(Comprehension)} For each of the following practices, say whether it is free,
requires a union bound, or is forbidden, with a one-sentence justification. (a)~Stopping
the gradient descent on the certificate when the certificate stops decreasing.
(b)~Choosing the scale $\sigma_0$ of a Gaussian prior among $5$ values after computing the
five certificates. (c)~Computing the certificate with $\delta=0.01$, reporting it with a $99\%$ claim if it is below a target value, and otherwise reporting the certificate computed with $\delta=0.05$ with a $95\%$ claim. (d)~Choosing the temperature of a
Gibbs posterior by cross-validation on the sample used in the bound.
\end{exercise}
\begin{exercise}
\textbf{(Comprehension)} True or false? Justify briefly. (a)~A self-bounding certificate
is an unbiased estimate of the true risk of the learned predictor. (b)~If the optimizer
of a self-bounding algorithm is stuck in a poor local minimum, the reported certificate
may be wrong. (c)~A test-set bound remains valid if the test set is also used to choose
among several predictors.
\end{exercise}
\begin{exercise}
\textbf{(Application)} In the model-selection example, the tandem certificates are computed
with $\delta/6$. For a forest with an OOB empirical joint error $\eS(Q)=0.05$,
$\KL(Q\|P)=2$ and $n_0=423$, compute the tandem certificate \eqref{eq:tnd_pb} with
$\delta=0.05$, with $\delta/6$, and with $\delta/100$. With a countable family and
$\pi_k=\frac6{\pi^2k^2}$, what is the extra cost, in nats, for the candidates $k=1$ and $k=10$?
\end{exercise}
\begin{exercise}
\textbf{(Application)} With $m=1000$ examples, the hold-out approach trains a forest on
$700$ examples and its vote errs on $30$ of the $300$ remaining ones. The self-bounding
approach trains a forest on the $1000$ examples and obtains an OOB empirical joint error of
$0.06$ with $\KL(Q\|P)=1.5$ and $n_0=370$. With $\delta=0.05$, compute the test-set
bound (Proposition~\ref{prop:testset}) and the tandem certificate. Below which value of
the empirical joint error would the tandem certificate beat the hold-out one? What is the
tandem certificate when $\eS(Q)=0$?
\end{exercise}
\begin{exercise}
\textbf{(Proof)} Prove Proposition~\ref{prop:union} in detail: write the failure event of each
bound, apply the union bound, and explain why the certificate $\min_k\mathcal B_k(Q_S,S,\pi_k\delta)$
is valid for any learning rule and any data-dependent choice of $k$. Check that
$\pi_k=\frac{6}{\pi^2k^2}$, $k\geq1$, is a valid choice of weights.
\end{exercise}
\begin{exercise}
\textbf{(Proof)} Let $\rho=\mathrm{softmax}(\theta)$ with $\theta\in\R^n$. Show that
$\frac{\partial\rho_s}{\partial\theta_t}=\rho_s(\ind{s=t}-\rho_t)$, deduce the formula
\eqref{eq:softmax_grad} for $\nabla_\theta\mathcal B$, and show that the components of
$\nabla_\theta\mathcal B$ sum to zero. Interpret this last fact.
\end{exercise}
\end{exercises}
